%% Section 6 -- Evaluation
%%
%% REWRITTEN 2026-08-29. This section previously reported no evaluation because
%% no implementation existed. It now reports reconstruction of synthetic
%% fragment pairs against known ground truth.
%%
%% THE SCOPE RULE THAT REPLACES THE OLD "DO NOT PAD" RULE: every number below
%% comes from a run that actually happened, and the section says what the
%% numbers do NOT establish as plainly as what they do. No real archive has been
%% reconstructed. Do not add a figure here that was not produced by executing
%% scripts/reconstruct.js, and do not describe synthetic ground truth in
%% language that would let a skimming reader take it for field data.

\section{Evaluation}
\label{sec:evaluation}

The reconstruction was run against the synthetic fragment pairs of
Section~\ref{sec:implementation:pairs}, whose ground truth is known because the
generator constructed it. This establishes that the specification is
implementable and that an implementation of it reproduces a known timeline. It
establishes nothing about real archives, which we return to at the end of the
section.

\subsection{Reconstruction against known ground truth}
\label{sec:evaluation:reconstruction}

The committed pair is generated with $\delta = 14$ days and $\epsilon = 7$ days
from a ground-truth timeline of 170 messages. The split gives $F_1$ 57 message
records and $F_2$ 133, with 20 messages falling in an overlap region and
therefore present in both fragments. Each fragment also opens with an
encryption-notice system line, which the two do not share.

Running Ingestion, Normalization, the declared fragment-scope step, and Synthesis
over both fragments yields 192 normalised records entering Synthesis and a
unified archive of 172. The 20 duplicated messages are collapsed to 20 canonical
records carrying provenance from both fragments, leaving 170 messages and the two
system lines. Nothing is quarantined and no conflict is reported.
$|\mathcal{U}|$ therefore equals the generator's ground truth exactly, and the
count of collapsed duplicates equals the constructed overlap exactly.

Conservation holds in the sense Invariant 1 requires: the 172 records in
$\mathcal{U}$, the 20 collapsed duplicates and the 0 quarantined records account
for all 192 records that entered. Invariants 1 and 2 are verified inside the
procedure rather than checked afterwards, so a run that violated either would
raise instead of returning an archive.

Regenerating the same pair with $\delta = \epsilon = 0$ gives the disjoint
configuration of the motivating incident. There, 172 records enter and 172
leave, with nothing collapsed and no conflict reported: Synthesis reduces to
ordered concatenation. The reduction happens without a branch for it anywhere in
the implementation, which is what Section~\ref{sec:carta:synthesis} claims should
happen and is the one claim in this paper that the disjoint case can test.

Two consecutive runs over identical input produce byte-identical
$(\mathcal{U}, \Omega)$, and supplying the fragments in the opposite order
produces the same archive. Invariant 3 is a requirement on implementations
rather than a theorem, and this is one implementation exhibiting it on one input,
not a proof.

\subsection{What the quarantine rule costs}
\label{sec:evaluation:quarantine}

The run above uses the fragment-scope step of
Section~\ref{sec:carta:normalization}. Running the same pair without it shows
what per-record judgement costs on its own, and the figure is large enough to be
worth reporting rather than smoothing over.

Judged per record, 81 of the 192 records are quarantined under
\texttt{AMBIGUOUS\_DATE\_FORMAT}: every record whose date has both leading
fields in the range 1 to 12 and which therefore parses under both conventions.
Only 111 records reach Synthesis, $|\mathcal{U}|$ falls to 103, and only 8 of the
20 overlapping messages are still deduplicated, because the other 12 have at
least one side quarantined. A reconstruction under the strict rule alone recovers
substantially less of the timeline than the fragments contain.

That is the specification behaving as designed rather than a defect. The point of
the rule is that a record whose date admits two readings does not silently
acquire one. The cost of the rule is visible here as a number, and the
fragment-scope step exists to recover from it while reporting that it did: in the
full run, 103 of the 172 canonical records were dated directly and 69 carry a
mark naming the inference, the fragment that supplied the evidence, and how many
records supplied it.

\subsection{The corpus's own boundary case}
\label{sec:evaluation:boundary}

One fixture in the five-file corpus, the group-chat fixture named in
Section~\ref{sec:implementation:inputs}, contains no record whose date admits a
single reading. All 25 of its records
quarantine, and the fragment-scope step declines to run on it, reporting that no
record in the fragment established an ordering. The other four fixtures resolve
completely once the step is applied.

This is the boundary the specification's per-record rule implies, appearing in an
artefact rather than in an argument. Where a fragment can settle its own
convention the declared step recovers the ambiguous records; where it cannot,
nothing in the fragment licenses a reading and the records stay quarantined. An
implementation that guessed would produce a complete-looking archive here, and a
reader would have no way to tell it apart from one that had evidence.

\subsection{What this does not establish}
\label{sec:evaluation:limits}

The ground truth these numbers are checked against was written by the same
project that wrote the implementation. That is the appropriate test for whether a
specification is implementable and self-consistent, and it is not evidence about
real archives. Specifically:

No real WhatsApp export has been reconstructed. The corpus is synthetic by
policy, since the project holds no permission to place real message content in a
public repository, and Section~\ref{sec:implementation:inputs} already limits the
nineteen edge-case categories to input forms this specification anticipates
rather than input forms observed in the field. A real archive may contain forms
the generator does not produce, and the implementation would quarantine or
mis-parse them.

No accuracy, precision or recall figure is reported, because the quantity those
would measure is agreement with a ground truth that here is a construction. The
counts above are conservation and identity checks, not measurements of a
classifier.

No performance measurement is reported. The largest input is 192 records, and
nothing about how the implementation behaves on an archive of a hundred thousand
has been tested.

The declared parameters of Section~\ref{sec:carta:conformance} were chosen once
and not varied. A different clustering rule or a wider deduplication window would
produce a different archive from the same fragments, which is what
Section~\ref{sec:algorithms:similarity} says openly; how much different has not
been measured.
